% SAMPLE-ONLY: \SampleGuideNote adds PDF guidance and is not a template command.
\section{まとめ}
\label{sec:summary}

本実験では、抵抗とコンデンサから構成される一次 RC 低域通過フィルタを作製し、
100 Hz から 10 kHz までの周波数応答を測定した。
入力振幅を $1.00\,\mathrm{V}$ に保ち、出力振幅と入出力波形の時間差から
利得と位相差を求めた。

低周波数側では利得が $0\,\mathrm{dB}$ に近く、入力信号がほぼ保たれた。
周波数の増加に伴って利得は低下し、位相差は $-90\,\mathrm{deg}$ に近づいた。
この変化は\cref{fig:gain-response,fig:dual-response}に示した理論曲線と同じ傾向であった。

遮断周波数は、利得の二点補間から 1.57 kHz、位相差の二点補間から 1.55 kHz、
利得曲線全体の当てはめから 1.58 kHz と推定された。
測定した素子値から計算した値は 1.58 kHz であり、曲線全体を用いた推定値と一致した。

主な定量結果を以下にまとめる。

% [DOC:LIST-ITEMIZE-04] Quantitative summary list.
\begin{itemize}[leftmargin=3em,itemsep=0.3em]
  \item $1.60\,\mathrm{kHz}$ における利得：$-3.10\,\mathrm{dB}$
  \item $1.60\,\mathrm{kHz}$ における位相差：$-46\,\mathrm{deg}$
  \item 曲線当てはめによる遮断周波数：$1.58\,\mathrm{kHz}$
  \item 公称素子値から求めた値との相対差：$0.63\,\%$
  \item 遮断周波数付近の利得標準偏差：$0.02$--$0.03\,\mathrm{dB}$
\end{itemize}

\SampleGuideNote[teal!70!black]{まとめの一覧に関する補足}{%
定量結果だけを短く再掲する場合にも、\texttt{itemize} を使用して本文と数値を分けられます。}

理論値と測定値の小さな差には、素子値の公差、測定周波数の間隔、
信号発生器の出力抵抗、プローブの入力容量、振幅と時間差の読取り分解能が含まれる。
三回の繰り返し測定では大きなばらつきは見られなかったが、系統的な差をさらに評価するには、
測定点を増やし、測定器の等価回路を理論計算へ取り入れる必要がある。

以上より、一次 RC 低域通過フィルタの振幅特性と位相特性を同じ測定から確認し、
素子値、表形式の測定結果、グラフおよび解析コードを対応させて遮断周波数を評価できた。
